\newpage
\section{Methodik und Forschungskonzept}

\subsection{Überblick und Forschungsdesign}

Im Folgenden werden die Methodik und das Forschungskonzept zur Beantwortung der Forschungsfrage dargestellt, welchen Einfluss der \ac{UI}-Kontext auf die Qualität \ac{LLM}-generierter Playwright-\ac{E2E}-Tests für \ac{OPT}-basierte WebGIS-Anwendungen hat. Die in der Einleitung formulierten Teilfragen werden dabei in das Untersuchungsdesign einbezogen.

Der \ac{UI}-Kontext ist die einzige unabhängige Variable. Alle anderen Faktoren wie Use Cases, LLM und dessen Konfiguration, OPT-Playwright-Skill sowie die manuell erstellten Referenztests werden konstant gehalten. Nur unter diesen Bedingungen sind die Qualitätsunterschiede der generierten Tests dem Kontext zurechenbar und werden nicht durch andere Einflüsse verfälscht. Dieses Vorgehen knüpft an das in Kapitel~2.1.1 beschriebene In-Context Learning an, bei dem der Kontext im Prompt die Ausgabe hauptsächlich bestimmt.

Von der unabhängigen Variable abzugrenzen sind die technischen Eigenschaften, die die Anwendung selbst zur Überprüfbarkeit bereitstellt. Während der \ac{UI}-Kontext beschreibt, was dem \ac{LLM} im Prompt mitgeteilt wird, bezeichnen Maßnahmen wie die Vergabe von data-testid-Attributen und der Zugriff auf das Kartenmodell über \texttt{\_\_openPioneerMap} (siehe Kapitel~5), was die Anwendung unabhängig vom Prompt für Tests bereitstellt. Diese Möglichkeiten bestehen bereits ab Stufe~1 und bleiben über alle vier Stufen unverändert. Variabel ist nur, ob und in welcher Form dem \ac{LLM} im Prompt mitgeteilt wird, dass sie existieren. Die Bereitstellung dieser technischen Möglichkeiten ist damit keine eigene Stufe des Experiments, sondern eine über alle Stufen geltende Voraussetzung, ohne die der \ac{UI}-Kontext nicht die einzige unabhängige Variable wäre.

Generiert werden die Tests aus den fachlichen Anforderungen des jeweiligen Use Case und nicht aus dem Quellcode der Anwendung. Die Struktur des Prompts wird dabei über alle vier Stufen konstant gehalten. Es ändert sich nur der enthaltene \ac{UI}-Kontext. Verschiedene Prompting-Strategien sind nicht Teil dieser Untersuchung. Da die Generierung nicht deterministisch erfolgt, wird jeder Use Case pro Stufe mehrfach generiert (N=50). Verglichen werden dadurch Verteilungen der Testqualität je nach Stufe, nicht einzelne Ausgaben.

Als kontrollierte Konstante im Prompt wird ein OPT-Playwright-Skill verwendet. Als OPT-Playwright-Skill wird im Rahmen dieser Arbeit eine feste, appunabhängige Fähigkeit bezeichnet, die dem \ac{LLM} bei jeder Generierung als Teil des Prompts mitgegeben wird. Sie legt allgemeine Umgangsweisen für den erzeugten Testcode fest, darunter die Verwendung von Playwright mit TypeScript, den Umgang mit \code|async|, \code|await|, \ac{OPT}-spezifische Besonderheiten sowie das erwartete Ausgabeformat. App-spezifische Selektoren enthält der Skill bewusst nicht, um eine Vermischung mit dem UI-Kontext zu vermeiden. Er wird unabhängig davon über alle vier Stufen unverändert angewendet. Die genaue Ausgestaltung des Skills wird in Kapitel~4.4 beschrieben.

Die vier Stufen bilden einen Verlauf von minimalem zu maximalem \ac{UI}-Kontext. Stufe~1 gibt dem \ac{LLM} keine Informationen über die Oberfläche, während die höheren Stufen den Kontext schrittweise erweitern und strukturieren. Abbildung~\ref{fig:ui_kontext_stufen} stellt die vier Ausprägungen mit steigendem Informationsgehalt dar.

\begin{figure}[H]
    \centering
    \caption[Ausprägungen der Variable UI-Kontext]{Die vier Ausprägungen der Variable UI-Kontext mit steigendem Informationsgehalt}
    \label{fig:ui_kontext_stufen}
    \includegraphics[width=0.6\textwidth]{ui_kontext_stufen}
\end{figure}

Wie ein einzelner Generierungsvorgang abläuft, zeigt Abbildung~\ref{fig:methodik_pipeline}. Use Case, OPT-Playwright-Skill und der \ac{UI}-Kontext der jeweiligen Stufe bilden zusammen den Prompt für das \ac{LLM}. Die daraus erzeugten Playwright-Tests durchlaufen ein zweiphasiges Evaluationsverfahren. In der ersten Phase werden sie gegen die laufende Demo-Anwendung ausgeführt und ihr Ausführungsergebnis klassifiziert. In der zweiten Phase bewertet ein zweites \ac{LLM} die Testqualität anhand fester Kriterien, wobei die manuellen Referenztests aus Kapitel~5.3 als Vergleichsmaßstab dienen. Beide Phasen werden in Kapitel~4.5 beschrieben.

\begin{figure}[H]
    \centering
    \caption[Generierungs- und Evaluationspipeline]{Generierungs- und Evaluationspipeline. Der UI-Kontext variiert als unabhängige Variable (Stufen 1-4), alle übrigen Eingaben sind konstant}
    \label{fig:methodik_pipeline}
    \includegraphics[width=1\textwidth]{methodik_pipeline}
\end{figure}

\subsection{Use Cases: Definition und Format}

Die Use Cases sind die fachliche Eingabe, aus der die Tests generiert werden. Sie sind über alle vier Stufen identisch und gehören damit zu den konstant gehaltenen Faktoren.

Jeder Use Case beschreibt eine Nutzerinteraktion mit der WebGIS-Anwendung auf fachlicher Ebene aus Sicht des Anwenders. Die Schritte enthalten bewusst keine Selektoren oder andere technischen Details der Oberfläche, damit das Wissen allein aus dem \ac{UI}-Kontext der jeweiligen Stufe kommt. Zudem ist es nicht sinnvoll, in einer Nutzerinteraktion technische Details der Anwendung festzuhalten, da der Anwender diese ebenfalls nicht kennt.

Die Use Cases werden im Markdown-Format erfasst. So sind sie menschenlesbar, für \acp{LLM} einfach zu verarbeiten und folgen einer einheitlichen Struktur, die für alle Stufen gleich verwendet wird. Jeder Use Case besteht aus Informationen zu \texttt{id}, \texttt{title}, \texttt{description}, \texttt{preconditions}, \texttt{steps}, \texttt{expected results} und \texttt{complexity}. Titel und Beschreibung geben an, welche Interaktion der Use Case abbildet. Die Vorbedingungen halten fest, was vor der Ausführung erfüllt sein muss. Die Schritte beschreiben die durchzuführende Interaktion des Anwenders und die erwarteten Ergebnisse legen fest, welcher Zustand danach eintreten soll.

Die Eigenschaft \texttt{complexity} gruppiert die Use Cases in drei Schwierigkeitsstufen: easy, medium, hard. Die Einstufung basiert nicht allein auf der Anzahl der Schritte, sondern vor allem auf der Art und Tiefe der Interaktionen. Einfache Use Cases prüfen einzelne Schaltflächen oder Sichtbarkeiten, komplexere fordern hingegen asynchrone Prozesse, Karteninteraktionen oder das Zusammenspiel mehrerer Komponenten. Die Abstufung stellt sicher, dass bei der Untersuchung verschiedene Komplexitäten typischer WebGIS-Interaktionen abgedeckt und nicht nur einfache Fälle betrachtet werden. Ausführliche Informationen zum Inhalt der Use Cases gibt es in Kapitel~5.2.

\subsection{Versuchsaufbau: Die vier Kontextstufen}

\subsubsection{Stufe 1: Baseline (kein UI-Kontext)}

In Stufe~1 bleibt der \ac{UI}-Kontext-Block des Prompts leer. Das \ac{LLM} erhält nur den Use Case und den OPT-Playwright-Skill, sonst nichts. Selektoren und Verifikationsstrategien muss es aus allgemeinem Wissen über Webanwendungen ableiten, z.\,B. durch das Raten plausibler Beschriftungen oder \ac{ARIA}-Rollen. Diese Stufe ist der Referenzpunkt des Experiments und soll zeigen, was ohne Informationen über die konkrete Oberfläche erreichbar ist. Die in Kapitel~4.1 beschriebenen Testbarkeitsmaßnahmen der Anwendungen existieren auch hier schon, dem \ac{LLM} wird ihre Existenz aber nicht mitgeteilt.

\subsubsection{Stufe 2: Automatischer Accessibility-Snapshot}

Stufe~2 erhält den \ac{UI}-Kontext automatisch aus der laufenden Anwendung. Ein Headless-Browser lädt die Demo-App und erfasst alle vergebenen data-testid-Werte und den Accessibility-Snapshot von Playwright, der Rollen, Namen und Zustände der Elemente als Baumstruktur ausgibt. Der Snapshot wird einmal vor dem Generierungslauf erfasst und für alle Use Cases identisch verwendet.

Der Ansatz benötigt keinen manuellen Pflegeaufwand und wäre damit in der Praxis ein günstiges Szenario. Erfasst wird allerdings nur der initiale Seitenzustand. Elemente, die erst nach Interaktionen erscheinen, fehlen im Snapshot. Auch das Kartenmodell taucht in diesem Kontext nicht auf, da \texttt{\_\_openPioneerMap} ein globales JavaScript-Objekt ist und weder im \ac{DOM} noch im Accessibility-Baum erscheint.

\subsubsection{Stufe 3: Automatisch generierte UI-Map}

Den gesamten Quellcode der Anwendung in den Prompt zu übernehmen ist wegen der begrenzten Kontextlänge des \ac{LLM} nicht praktikabel. Stufe~3 verwendet deshalb eine automatisch aus dem Quellcode generierte \ac{UI}-Map. Die regexbasierte Extraktion durchsucht alle \texttt{.ts}- und \texttt{.tsx}-Dateien der Anwendung und erfasst data-testids, Beziehungen zwischen Komponenten, \ac{UI}-Zustände, Sichtbarkeitsbedingungen und Interaktionsabhängigkeiten. Informationen zur Sichtbarkeit der Layer und Hintergrundkarten werden zusätzlich aus der Service-Konfiguration (\texttt{services.ts}) gelesen. Das Ergebnis wird dann als Markdown-Tabelle gespeichert und in den Prompt übernommen. Zusätzlich enthält der Prompt erstmals die Datei \texttt{map-model-helpers.ts}. Deren Funktionen helfen beim Zugriff auf das Kartenmodell und geben dem \ac{LLM} damit ein Mittel, den Kartenzustand in Assertions zu prüfen.

Der Kontext in dieser Stufe stammt aus dem Quellcode statt aus der laufenden Seite. Er enthält deshalb auch Elemente, die im initialen Zustand nicht sichtbar sind. Der Ansatz bleibt automatisiert, setzt aber Zugriff auf den Quellcode voraus, und seine Qualität hängt davon ab, wie zuverlässig die regexbasierte Extraktion die relevanten Informationen findet.

\subsubsection{Stufe 4: Manuell erstellte UI-Map}

Stufe~4 folgt dem Aufbau von Stufe~3, ersetzt die generierte \ac{UI}-Map allerdings durch eine manuell erstellte Fassung. Sie wurde mit vollständiger Kenntnis der Anwendung erfasst und enthält alle Informationen zu Selektoren sowie Hinweise zu Interaktionen und Zuständen einzelner Komponenten. Anders als die Tabelle aus Stufe~3 bildet sie die Komponenten hierarchisch ab und modelliert dynamische Sichtbarkeiten mit expliziten Bedingungen. Damit diese Beziehungen effizient dargestellt werden können, werden die Informationen im \ac{JSON}-Format gespeichert. Die Map-Model-Helfer sind identisch zu Stufe~3.

Diese Stufe bildet die theoretische Obergrenze des untersuchten Ansatzes ab. Sie soll zeigen, welche Testqualität erreichbar ist, wenn der Kontext vollständig und geprüft vorliegt. Dem steht allerdings der Aufwand für die Erstellung und Pflege gegenüber. Jede Änderung an der Oberfläche muss in der \ac{UI}-Map nachgezogen werden, sonst enthält sie veraltete Selektoren oder falsche Informationen zu Sichtbarkeiten und Interaktionen. In der Evaluation wird sich zeigen, ob der Qualitätsgewinn diesen Aufwand im Vergleich zu den Testergebnissen der anderen Stufen rechtfertigt.

\subsubsection{Bonus-Stufe 5: Self-Improvement-Loop}

Ergänzend zum Vergleich der vier Stufen wird in der Bonus-Stufe~5 eine Selbstverbesserungsschleife untersucht. Sie dient als Erweiterung der anderen Stufen und verfolgt einen etwas anderen Ansatz. Statt den statischen \ac{UI}-Kontext im Prompt weiter zu verbessern, erhält das \ac{LLM} Rückmeldung aus der Ausführung seines eigenen Tests, um diesen weiter zu optimieren.

Als Ausgangspunkt dient der Kontext aus dem automatischen Snapshot des initialen Seitenzustands aus Stufe~2, ergänzt um einen Screenshot der Anwendung und den Map-Model-Helfer aus Stufe~3 und 4. Den Ablauf der Schleife zeigt Abbildung~\ref{fig:stufe5_loop}. Zuerst wird ähnlich wie bei Stufe~2 aus dem initialen Kontext ein Test generiert und mit Playwright direkt gegen die laufende Anwendung ausgeführt. Besteht er, endet die Schleife. Schlägt er fehl, wird der nächste Prompt um den fehlgeschlagenen Testcode, die Playwright-Fehlermeldung und den Failure-Snapshot inklusive Screenshot ergänzt. Das Modell erhält damit die Anweisung, den vorhandenen Test zu korrigieren, statt einen neuen zu erzeugen. Nach spätestens 10 Iterationen bricht die Schleife ab.

\begin{figure}[H]
    \centering
    \caption[Self-Improvement-Loop (Stufe 5)]{Self-Improvement-Loop der Bonus-Stufe 5: Fehlermeldung und Failure-Snapshot inkl. Screenshot fließen in den nächsten Generierungsprompt zurück (max. 10 Iterationen)}
    \label{fig:stufe5_loop}
    \includegraphics[width=0.8\textwidth]{stufe_5_loop}
\end{figure}

Der Failure-Snapshot hält den Zustand der Anwendung zum Zeitpunkt des Fehlschlags fest und besteht wie in Stufe~2 aus Accessibility-Snapshot und data-testids. Der Unterschied liegt im Zeitpunkt der Aufnahme. Hier wird nicht der initiale Zustand, sondern die Seite nach den bereits ausgeführten Testschritten erfasst. Das \ac{LLM} sieht damit, wie die Oberfläche aussah, als der Test scheiterte. Zusätzlich erhält das \ac{LLM} in der initialen und jeder weiteren Iteration einen Screenshot der Anwendung als Bildeingabe. Dadurch bekommt es auch erstmals visuelle Informationen über den gerenderten Kartenzustand und sieht optisch, welche Elemente sichtbar sind. Als Pass-Kriterium dient der Screenshot nicht. Die Schleife endet nur über das Ausführungsergebnis des Tests. Schlägt der Test beispielsweise fehl, obwohl der Layer im Screenshot sichtbar ist, liegt der Fehler in der Prüflogik des Tests, und die Iteration wird fortgesetzt. Der Screenshot ist nur ein Hilfsmittel für die Korrektur.

Jede Iteration verwendet einen eigenständigen Prompt ohne Konversationshistorie. Alle Informationen aus dem vorherigen Durchlauf werden ausschließlich über den vorherigen Testcode, die Fehlermeldung und den Snapshot übergeben. Für die Aufnahme des Failure-Snapshots wird eine Playwright-Fixture benötigt, die in einer Ausführungskopie des Tests eingebunden wird. Die vom \ac{LLM} generierte Originaldatei bleibt dabei unverändert und ist Grundlage der Bewertung.

\subsubsection{Erwartungen}

\subsection{LLM-Konfiguration (Modell, Temperatur, Prompt-Struktur)}

Für alle Generierungen wird das Qwen3.6-35B-A3B als \ac{LLM} verwendet. Es ist ein offenes Sprachmodell des Qwen-Teams von Alibaba und wird in der FP8-quantisierten Variante eingesetzt. Das Modell verwendet eine Sparse-Mixture-of-Experts-Architektur mit 35 Milliarden Parametern insgesamt, von denen pro Token nur etwa 3 Milliarden aktiv sind. Zudem ist es multimodal, verarbeitet also neben Text auch Bilder. Insgesamt unterstützt es eine Kontextlänge von 262.144 Token (\cite{hugging_face_qwenqwen36-35b-a3b_2026, qwen_team_alibaba_qwen36-35b-a3b_2026}).

Das Modell wird auf einem Spark im Netzwerk der con terra GmbH betrieben und dort über LM~Studio bereitgestellt. Die Generierungsskripte greifen über die OpenAI-kompatible \ac{REST}-\ac{API} von LM~Studio darauf zu.

Die Temperatur ist auf 0.6 gesetzt, da dieser Wert vom Hersteller für präzise Coding-Aufgaben empfohlen wird (\cite{hugging_face_qwenqwen36-35b-a3b_2026}). Qwen3.6 arbeitet standardmäßig im Thinking-Mode, der auch hier verwendet wird. Die maximale Ausgabelänge beträgt 64k Token, damit Reasoning und Testcode zusammen ausreichend Platz haben. Diese Konfiguration ist über alle Stufen und Use Cases identisch. Da die Generierung bei einer Temperatur über null nicht deterministisch ist, wird der vollständige Generierungsdurchlauf pro Stufe 50-mal wiederholt. Jeder Durchlauf erzeugt alle zehn Use Cases in einem eigenen Verzeichnis, sodass pro Stufe 500 Tests generiert werden.

Der Prompt besteht aus zwei Teilen. Der User-Prompt enthält die Anweisung, aus dem Use Case mit dem jeweiligen \ac{UI}-Kontext einen Playwright-\ac{E2E}-Test zu generieren, den Use Case selbst, die \ac{URL} der Anwendung und den \ac{UI}-Kontext-Block der jeweiligen Stufe. Zwischen den Stufen variiert ausschließlich der \ac{UI}-Kontextblock. Im System-Prompt wird der OPT-Playwright-Skill hinzugefügt.

Der Skill legt fest, wie der erzeugte Testcode auszusehen hat. Er schreibt die Verwendung von Playwright mit TypeScript vor, verlangt die durchgängige Nutzung von \code|async| und \code|await|, priorisiert \code|getByTestId| gegenüber anderen Locator-Strategien und definiert das Ausgabeformat. Die Ausgabe soll eine einzelne, vollständige Testdatei ohne erläuternden Text sein. Außerdem enthält der Skill weitere technische Vorgaben, etwa zum Umgang mit Lokatoren, Assertions, Wartebedingungen und Chakra-UI-Elementen. Der Skill ist so aufgebaut, dass er auch für andere WebGIS-Anwendungen wiederverwendet werden kann, die mit Open Pioneer Trails erstellt wurden. Informationen über App-spezifische Selektoren oder Elemente enthält er daher bewusst nicht.

\subsection{Evaluationsverfahren}

\subsubsection{Phase 1: Deterministische Ausführung}

\subsubsection{Phase 2: LLM-as-a-judge}
